\section{Contraction by multiplicative cotangent systems}
\label{g:multiplicative}

The contraction in Section~\ref{g:cf:contraction} chooses systems of
zero degree for one supporting movable class. There is another useful
choice: close all cotangent-twisted systems under multiplication.
This avoids fixing that class, while retaining the same decisive
verification of generic adjoint rank after contraction. We state the
construction with its prepared data so that its difference from the
zero-degree method is explicit.

\begin{proposition}\label{g:multiplicative-contraction}
Let $X$ be smooth projective and rationally connected, let
$L=-K_X$ be smoothly semipositive, and let an algebraic torus $T$ act
with the natural linearization. Let
$\pi:Z\to X$, $g:Z\to Y$ resolve a dominant invariant rational map,
with $Y$ smooth projective of positive dimension and
$\mathbb C(Y)$ relatively algebraically closed in $\mathbb C(X)$.
Write $E=K_{Z/X}$. Suppose a nonzero invariant rational section $s$ of
$\ell\pi^*L$, $\ell>0$, and rational divisors $N$ on $Y$ and $D$ on
$Z$ satisfy
\[
 \ell^{-1}\operatorname{div}_Z(s)=g^*N+D,\qquad D\geq0.
\]
Assume that $N$ has a semipositive metric with locally bounded weights.
Suppose there is a rational simple normal crossings divisor
$\Delta=\Delta^+-\Delta^-$ on $Y$, with every coefficient of
$\Delta^+$ less than one, such that $-K_Y-\Delta$ is nef.
Let $B$ be a reduced effective divisor on $Y$ such that
\begin{enumerate}[label=(\roman*)]
\item $\operatorname{Supp}\Delta^-\subset\operatorname{Supp}B$;
\item $\pi_*(D-tg^*B)\geq0$ for every sufficiently small rational $t>0$;
\item above every base prime outside $B$, some prime dominating it has
coefficient zero in both $D$ and $E$.
\end{enumerate}
Finally suppose the invariant generic spaces of
$E+k\pi^*L$ have rank one for $k=0$ and every positive multiple of
one fixed integer.
Then either $H^0(X,rL)^T\ne0$ for some $r>0$, or a rational map
$Y\dashrightarrow Y'$ with $\dim Y'<\dim Y$ gives a composite
fibration retaining a nonzero invariant generic section in a positive
multiple of $L$ and invariant adjoint rank one at zero and in every
sufficiently divisible positive degree.
\end{proposition}

These prepared data are supplied by vertical minima and the toroidal
logarithmic density calculation of
Lemma~\ref{g:cf:metric-conclusions}; no strict-curvature patch is
required. The signed slope statement needed in the proof is
Proposition~\ref{g:bt:signed-slope}, proved for these signed boundaries.
In particular the argument does not apply an effective-boundary
positivity theorem directly to $\Delta$.

\begin{proof}
Suppose first that $N+tB$ is big for some sufficiently small rational
$t>0$. Choose a nonzero section of a divisible multiple. Pull it back,
use the corresponding power of $s$, and push its divisor to $X$.
Condition (ii) leaves an effective divisor. All rational coefficients
can be cleared by the power, so this is a regular nonzero invariant
section of a positive multiple of $L$. We may therefore assume that
$N+tB$ is not big for all small positive rational $t$.

For every line bundle $M$ admitting a nonzero map
$M\to(\Omega_Y^1)^{\otimes p}$,
Proposition~\ref{g:bt:signed-slope} gives
\begin{equation}\label{g:multiplicative-cotangent}
 p\Delta^- -M\quad\text{pseudoeffective}.
\end{equation}
The case $p=0$ follows directly from $M\to\mathcal O_Y$.
Choose an integer $b>0$ clearing all divisor denominators, divisible by
$\ell$, and compatible with the generic-rank hypothesis. If $s_{bN}$
denotes the canonical rational section of $bN$, the normalized section
\[
 \sigma=\frac{s^{b/\ell}}{g^*s_{bN}}
 \quad\text{of }b\pi^*L-g^*(bN)
 \quad\text{has divisor }bD\geq0.
\]
It is regular and invariant. Consider the
line bundles
\begin{equation}\label{g:multiplicative-collection}
 Q=kN+aB+M,
 \quad k\in b\mathbb Z_{\geq0},\quad a\in\mathbb Z_{\geq0},
 \quad M\longrightarrow(\Omega_Y^1)^{\otimes p},\quad p\geq0,
\end{equation}
where the last map is nonzero. This collection is closed under tensor
product: tensoring two nonzero line maps gives a nonzero map at the
generic point and hence a nonzero global map into the corresponding
tensor power.

No member of \eqref{g:multiplicative-collection} is big. Indeed,
\eqref{g:multiplicative-cotangent} would make
$kN+aB+p\Delta^-$ big if $Q$ were big. Choose $a'\geq a$ with
$a'B-aB-p\Delta^-\geq0$. Then $kN+a'B$ would be big. For any prescribed
small $t>0$, choose $\delta>0$ so that $\delta k<1$ and
$\delta a'\leq t$. Since $N$ and $B$ are pseudoeffective,
\[
 N+tB=\delta(kN+a'B)+(1-\delta k)N+(t-\delta a')B
\]
would be big, contrary to the assumption.

Choose a nonempty system in a member $Q_0$ of
\eqref{g:multiplicative-collection} whose image has maximal dimension.
The constant system shows that such systems exist. Resolve its map and
pass to the connected-fiber base to obtain
\[
 \eta:\widetilde Y\to Y,\qquad q:\widetilde Y\to Y'.
\]
We can choose $Y'$ normal projective; its smooth birational model may
be used later. Its dimension is strictly smaller than $\dim Y$, since
$Q_0$ is not big. A suitable power of the chosen system supplies an
ample line $A$ on the normal Stein base with
$\eta^*Q_0-q^*A$ effective, after replacing $Q_0$ by that power.
Every ratio from every system in
\eqref{g:multiplicative-collection} belongs to $\mathbb C(Y')$.
Otherwise the product with the chosen system would increase the image
dimension; relative algebraic closure removes the remaining algebraic
extension.

For the generic fiber $F$ of $q$ this implies
\begin{equation}\label{g:multiplicative-rigidity}
 h^0(F,\eta^*Q|_F)\leq1
 \quad\text{for every }Q\text{ in }
 \eqref{g:multiplicative-collection}.
\end{equation}
To check the assertion for generic sections, spread two putative
independent sections after a sufficiently ample base twist $q^*(nA)$.
The effective difference $\eta^*(nQ_0)-q^*(nA)$ puts them in
$\eta^*(Q+nQ_0)$. Their pushdowns are sections on the smooth $Y$,
and their ratio contradicts the preceding field statement.

Because $Y$ is a rational image of the rationally connected $X$, it
has no positive-degree holomorphic forms. Apply
Lemma~\ref{g:determinant} to $bN$. We obtain one cotangent line $M_0$
and effective divisors $D_i\sim M_0+m_iN$ for an unbounded increasing
sequence $m_i\in b\mathbb Z_{>0}$. Their restrictions to $F$ are
nonzero section divisors. By \eqref{g:multiplicative-rigidity}, the
products realizing any three-index linear relation are proportional
on $F$. Therefore each horizontal multiplicity is affine in $m_i$.
All these multiplicities are nonnegative at unbounded exponents, so
the slope is nonnegative. A difference $D_j-D_i$, $j>i$, consequently
gives a nonzero section of $(m_j-m_i)\eta^*N$ on $F$.

Write $\delta=m_j-m_i$, and choose $\gamma\in\mathbb C(Y)^*$ with
\[
 D_j-D_i=\delta N+\operatorname{div}_Y(\gamma).
\]
The integer $\delta$ is divisible by $b$ and hence by $\ell$.
Resolve the composite map from $X$ to $Y'$ by a smooth model $Z'$
dominating $Z$ and $\widetilde Y$; write $j:Z'\to Y'$ and continue
to write $g:Z'\to Y$, $\pi:Z'\to X$, with $D$ pulled back from $Z$.
The rational section
\[
 \tau=s^{\delta/\ell}g^*\gamma\quad\text{of }\delta\pi^*L
 \quad\text{satisfies}\quad
 \operatorname{div}(\tau)=\delta D+g^*(D_j-D_i).
\]
The first summand is effective, and the second is effective on the
complete generic fiber over $Y'$, as just proved. Thus $\tau$ is a
nonzero invariant generic section. The composite generic fiber is geometrically
connected, since both successive function-field inclusions are
relatively algebraically closed in characteristic zero.

It remains to restore the adjoint ranks on this larger fiber.
Suppose two invariant sections of $K_{Z'/X}+k\pi^*L$ are independent
over $\mathbb C(Y')$, for a sufficiently divisible $k\geq0$.
Clear their poles over $Y'$ by $j^*(nA)$, and pull back the effective
difference $\eta^*(nQ_0)-q^*(nA)$ to regard them as sections of
\[
 K_{Z'/X}+k\pi^*L+g^*(nQ_0).
\]
Over $\mathbb C(Y)$ the old rank is one. Divide the two sections by
the Jacobian section and $\sigma^{k/b}$, using the pullback of the
normalized section defined above. The quotients descend to rational
sections $v_1,v_2$ of $kN+nQ_0$ on $Y$.
For a prime $P$ outside $B$, condition (iii) supplies a prime $R$
above it whose strict transform has order zero in both $D$ and
$K_{Z'/X}$. If $e_R=\operatorname{ord}_R g^*P$, regularity of the
two sections upstairs gives
\[
 0\leq\operatorname{ord}_R(s_{K_{Z'/X}}\sigma^{k/b}g^*v_i)
      =e_R\operatorname{ord}_P(v_i),\qquad i=1,2.
\]
Thus neither $v_i$ has a pole outside $B$. A sufficiently large $a$ therefore makes
both sections regular in $kN+nQ_0+aB$, which again belongs to
\eqref{g:multiplicative-collection}. Their ratio lies in
$\mathbb C(Y')$, a contradiction.

The new generic section, multiplied by the Jacobian section, supplies
a nonzero adjoint section at every positive multiple of its exponent;
the Jacobian supplies one at zero. Combining these lower bounds with
the upper bound proves all required ranks after enlarging the fixed
divisibility integer. Finally replace the normal Stein base $Y'$ by a
smooth projective model and resolve the composite map. Its function
field is unchanged, and generic sections and Jacobian-twisted ranks
are unchanged by this birational replacement, as in
Section~\ref{g:cf:generic-conditions}. This completes the contraction.
\end{proof}
